% 제출본 전용 — 9장 끝에 붙는다 (9.8). 공학인증 구현·시험평가 표.

\subsection{설계 구성요소 연관성 --- 시험 및 평가}

9장이 통째로 시험 및 평가에 해당한다. 공학 설계의 구성요소로 정리하면 다음과
같다.

\begin{table}[htbp]
  \centering
  \caption{설계 구성요소 연관성 --- 시험 및 평가}
  \label{tab:v2-test}
  \begin{tabular}{llcp{9.0cm}}
    \toprule
    \textbf{설계 구성요소} & \textbf{구분} & \textbf{연관성} & \textbf{내용} \\
    \midrule
    \multirow{2}{*}{설계 구성요소}
      & 구현 & O &
        수집부터 보고서의 표·그림 생성까지 \textbf{\pipeSteps 단계
        파이프라인} 하나로 돈다(3.5절). 자료 \pipeData · 정책 \pipePolicy ·
        분석 \pipeAnalysis · 재현 \pipeReplicate · 시각화 \pipeViz ·
        문서 \pipeDocs 단계이며, 뒤 단계가 앞 단계 산출물을 읽으므로
        하나라도 실패하면 거기서 멈춘다 \\
    \cmidrule(l){2-4}
      & 시험평가 & O &
        결과를 \textbf{스스로 깨뜨려 보는 검정 일곱}(9장) --- 위약 ·
        합성대조군 · 오염 제외 · 처치 역전 · 사양 곡선 \pipeAnalysis 축 ·
        외적 타당성 · 취소거래. 여기에 \textbf{자동 점검 넷}(표
        \ref{tab:checks})이 매 실행마다 걸리고, 새 폴더 새 가상환경에서
        돌려 결과 파일 25개가 \textbf{바이트 단위로 같은 것}을 확인했다
        (3.6절) \\
    \bottomrule
  \end{tabular}
\end{table}

\textbf{시험의 성격이 일반적인 소프트웨어와 다르다.} 기능이 명세대로 동작하는지를
보는 것이 아니라, \textbf{추정된 수치가 설계 선택의 산물이 아닌지}를 본다. 그래서
``통과/실패''가 아니라 ``가릴 수 있다/없다''로 결과가 나온다. 9.4절의 해제 효과와
10.4절의 지속 기간이 \textbf{가릴 수 없다}로 끝나는 것이 그 예이며, 이 역시
시험의 결과다.

\textbf{사용자설명서에 해당하는 것은 재현 절차다}(부록 C). 이 작품의 사용자는
화면을 쓰는 사람이 아니라 \textbf{결과를 다시 돌려 보려는 사람}이므로, 설치 ---
스냅샷 준비 --- 실행 --- 대조까지를 명령 단위로 적었다. 실제로 그 절차를 따라
클린 환경에서 돌려 보았고, 그 과정에서 의존 목록의 구멍 둘을 찾았다.
